\documentclass[12pt]{article}


\usepackage{amsmath}
\usepackage{amsthm}
\usepackage{amssymb}
\usepackage[T1]{fontenc}
\usepackage{newtxmath,newtxtext}
\theoremstyle{definition}%
\newtheorem{theorem}{Theorem}[section]
\newtheorem{definition}{Definition}[section]
\newtheorem{example}[theorem]{Example}

\begin{document}

\section{events}
\begin{definition}[sample space]
	The set of all possible outcomes of an experiment is called the $\textbf{sample space}$ and is denoted by $\Omega$

\end{definition}


\begin{example}
	Head or Tails has sample space $\Omega = \{H, T\}$,

	
\end{example}

\begin{definition}[Field]
	A field $\mathcal{F}$ is a subcollection of the collection of all subsets of $\Omega$. 
	\begin{align*}
		\mathcal{F} \subseteq \{A \colon A \subseteq \Omega\}
	\end{align*}
	This subcollection needs to satisfy:
	\begin{enumerate}
		\item 	if $A, B \in \mathcal{F}$, then $A \cup B \in \mathcal{F}$ and $A \cap B \in \mathcal{F}$

		\item if $A \in \mathcal{F}$ then $A^c \in \mathcal{F}$
		\item the empty set $\emptyset$ belongs to $\mathcal{F}$
	\end{enumerate}
	It follows from the properties of a field $\mathcal{F}$ that
	\begin{align*}
		\text{ if } A_1,A_2,\ldots, A_n \in \mathcal{F} \text{ then } \bigcup_{i = 1}^{n}A_i \in \mathcal{F}
	\end{align*}
\end{definition}
We note that $\omega$ is notation used for a member of $\Omega$
\end{document}


